\documentclass[11pt]{article}

\usepackage{amsmath,amssymb}
\usepackage{xurl}
\usepackage[hidelinks]{hyperref}
\setlength{\emergencystretch}{2em}

% Shared commands for notes in docs/paper-gaps/.
%
% Two halves.  The link commands below are project identity: OWNER/REPO is the
% placeholder that scripts/bootstrap_project.py replaces with this project's
% GitHub slug and Pages site.  Everything after them is NOTATION, and notation
% belongs to the source paper: the definitions here are an example set, and a
% project replaces them with the macros of its own paper's style file.  The one
% rule that never changes: a command defined here must mean exactly what the
% source paper's macro means, or a note silently says something else.

\newcommand{\Lean}{\textsc{Lean}}
\newcommand{\leanid}[1]{\textsf{#1}}
\newcommand{\ghissue}[1]{\href{https://github.com/OWNER/REPO/issues/#1}{GitHub issue \##1}}
\newcommand{\ghpr}[1]{\href{https://github.com/OWNER/REPO/pull/#1}{PR \##1}}
% Local issue tracker (issues/NNNN-*.md in this repository); no remote URL.
\newcommand{\localissue}[1]{issue \##1 (\path{issues/})}
\newcommand{\doclink}[1]{\href{https://OWNER.github.io/REPO/docs/#1.html}{\path{#1}}}

% --- Notation: replace with the source paper's own macros --------------------
% \F, \N, \C, \R, \tr, \ot, \eps, \E, \bx, \bu, \bv, \by

\newcommand{\F}{\mathbb{F}}
\newcommand{\N}{\mathbb{N}}
\newcommand{\C}{\mathbb{C}}
\newcommand{\R}{\mathbb{R}}
\newcommand{\tr}{\mathrm{tr}}
\newcommand{\eps}{\epsilon}
\newcommand{\ot}{\otimes}
\newcommand{\E}{\mathop{\bf E\/}}

% bold random variables
\newcommand{\bu}{\boldsymbol{u}}
\newcommand{\bv}{\boldsymbol{v}}
\newcommand{\bx}{{\boldsymbol{x}}}
\newcommand{\by}{\boldsymbol{y}}

% T1 so literal < and > in \texttt placeholders typeset as themselves
\usepackage[T1]{fontenc}

% bra-ket
\usepackage{braket}

% --- Example structured spaces ---
\newcommand{\polyfunc}[3]{\mathcal{P}(#1, #2, #3)}
\newcommand{\polymeas}[3]{\mathrm{PolyMeas}(#1, #2, #3)}
\newcommand{\polysub}[3]{\mathrm{PolySub}(#1, #2, #3)}

% Approximation relations (paper uses \approx_\eps and \simeq_\zeta)
\newcommand{\approxeps}{\approx_{\varepsilon}}
\newcommand{\simgeq}{\simeq_{\zeta}}


\renewcommand{\ghpr}[1]{%
  \href{https://github.com/Dengnifer/MIPStarRE-A/pull/#1}{PR \##1}}

\title{A Direct Construction of the Combined Point Measurements}
\author{}
\date{2026-09-09}

\begin{document}
\maketitle

\section{At a glance}

\begin{itemize}
  \item \textbf{Difficulty: moderate.}  The Pauli basis analysis combines
  binary refinements of the expanded point measurements, then uses quantum
  linearity to recover field-valued outcomes.  We instead combine the
  field-valued measurements directly.  A finite Parseval identity transfers
  the commutation estimate without a field-size loss, and an exact sandwich
  identity supplies the consistency needed for orthonormalization.
  \item \textbf{Estimated weight:} the completed construction comprises
  roughly two thousand lines of supporting formal proofs.
  \item \textbf{Mathlib/project split:} approximately \textbf{20\% Mathlib}
  and \textbf{80\% project-local}.  The project-local portion consists of the
  finite Fourier identities, register placements, and state-dependent distance
  estimates.
  \item \textbf{Key Mathlib inputs:} finite sums, complex inner products,
  matrix adjoints, and real powers, together with the existing projective
  rounding theorem.  No new stability theorem is assumed.
\end{itemize}

This proves the combined-point conclusion on the original expanded local
spaces without using the common-ancilla form of quantum linearity.\footnote{The construction is in
\path{MIPStarRE/QPBT/Combining/Points.lean} and its supporting modules, reviewed
in \ghpr{424}.  The independently proved common-ancilla theorem and the
distinct padding limitation are documented in
\path{docs/paper-gaps/qpbt_linearity-theorem-quotation.tex}.}

\section{Key theorem forms}

Fix admissible Pauli basis parameters, with $q=2^t$ and $M=2^m$, and let a
projective strategy pass the test with probability at least $1-\eps$.
Write $V=(\C^q)^{\ot M}$ for one Pauli register and
$\hat\psi=\psi_{AB}\ot\mathrm{EPR}_{A'A''}\ot\mathrm{EPR}_{B'B''}$
for its normalized expanded state.  For points $x,z\in\F_q^m$, write
$X_a^x=\hat M^{(\mathrm{Point},X),x}_a$ and
$Z_b^z=\hat M^{(\mathrm{Point},Z),z}_b$ for the projective point effects
on either player's space $H\ot V$.  A subscript $p$ places an operator on
one of $AA',BA'',BB',AB''$.  For operator families $F,G$ with outcomes $(a,b)$,
define
\[
 d^2(F,G)=\E_{x,z}\sum_{a,b}
   \bigl\|(F^{x,z}_{a,b}-G^{x,z}_{a,b})\hat\psi\bigr\|^2,
\]
where $x,z$ are uniform.  The combined-point lemma asserts that there are
projective measurements $Q^{x,z}$ on each original space $H\ot V$ such that
\[
 \begin{aligned}
 d^2(Q_p,Q_{p'})&\leq\delta_Q(\eps),\\
 d^2(Q_p,(XZ)_{p'})&\leq\delta_Q(\eps),\\
 d^2(Q_p,(ZX)_{p'})&\leq\delta_Q(\eps)
 \end{aligned}
 \qquad\text{for every directed opposite pair }(p,p').
\]
Here $(XZ)_{a,b}=X_a^xZ_b^z$, $(ZX)_{a,b}=Z_b^zX_a^x$, and the opposite
pairs are $AA',BA''$ and $BB',AB''$, in both directions.  We prove these
bounds with $\delta_Q(\eps)=K\eps^{1/8}$ for one universal constant $K$.
The hypotheses and the three conclusions are unchanged; the replacement is
in the proof, not in the statement.

\section{The source argument}

The combined-point lemma in the subsection ``Combining the $X$ and $Z$
measurements'' of \cite{Ji2020MIPStar} asserts the three conclusions above
with polynomially small error.\footnote{See
\path{references/qpbt-paper/14_analysis_of_the_pauli_basis_test.tex},
lines~689--709, label \path{lem:qld-4-10}.}
Its proof first constructs binary-outcome sandwiches, proves consistency by
a Cauchy--Schwarz calculation, and orthonormalizes them.  Approximate
linearity then permits a second construction with outcomes in
$\F_q^2$.\footnote{The binary construction and its consistency argument are
at lines~711--790 of the same source, including \path{eq:qld-r-2} and
\path{eq:qld-rw-self-cons-1} through \path{eq:qld-rw-self-cons-4}.
Lines~825--832 apply \path{thm:linearity} and invoke zero-state ancillas.}

The linearity theorem returns an auxiliary space for each point pair.  Its
common-ancilla form chooses one finite ancillary space and vector before
quantifying over the finite-dimensional observable families.  Thus one space
serves every point pair.\footnote{See
\leanid{MIPStarRE.QPBT.exists\_exactly\_linear\_observables\_commonAncilla}
in \path{MIPStarRE/QPBT/Combining/Linearity.lean}.}
The fixed expanded-state setting still does not supply the zero-state padding
assumed in the source.  The direct proof below constructs the measurements
on the original spaces without requiring that padding.

\section{Parseval and the sandwich defect}

Fix $x,z$ and one placement, suppressing these indices for the moment.
For the trace character $\chi(u)=(-1)^{\tr(u)}$, the expanded observables
$U_X^r,U_Z^s$ give the Fourier expressions
$X_a=q^{-1}\sum_r\chi(ar)U_X^r$ and
$Z_b=q^{-1}\sum_s\chi(bs)U_Z^s$.
With $[T,U]=TU-UT$, character orthogonality therefore gives
\[
 \sum_{a,b}\|[X_a,Z_b]\hat\psi\|^2
   =\frac1{q^2}\sum_{r,s}\|[U_X^r,U_Z^s]\hat\psi\|^2.
\]
Indeed, the coefficient matrix with rows $(a,b)$ and columns $(r,s)$ is
$c_{(a,b),(r,s)}=\chi(ar)\chi(bs)$, and its columns satisfy
\[
 \sum_{a,b}\overline{c_{(a,b),(r,s)}}c_{(a,b),(r',s')}
   =q^2\mathbf{1}_{(r,s)=(r',s')}.
\]
Expanding the squared norms as inner products cancels all unequal column
indices.  The squared Fourier prefactor $q^{-4}$ leaves $q^{-2}$, precisely
the uniform average above.  Thus the expanded-observable commutation bound
underlying the source's binary-refinement estimate yields a field-valued
bound $O(\sqrt\eps)$, with no dependence on $q$.\footnote{The Fourier
definition is at lines~384--418 of the source; the commutation estimate is
in \path{lem:qld-comm-cons}, lines~466--505.  The formal identities are
\path{sum_norm_sum_smul_sq_of_orthogonal} and
\path{ProjectiveSetting.sum_norm_place_expPointOp_commutator_sq};
\path{ProjectiveSetting.expPoint_comm} derives the bound in
\path{MIPStarRE/QPBT/Combining/Points/Commutation.lean}.}

Define the field-valued sandwich by $R_{a,b}=Z_bX_aZ_b$.
It is positive, since $R_{a,b}=(X_aZ_b)^\dagger(X_aZ_b)$, and its effects
sum to the identity by projectivity and completeness.  Moreover,
$R_{a,b}-Z_bX_a=Z_b[X_a,Z_b]$.  Since $Z_b$ is a contraction, the same
commutation bound controls $d^2(R,ZX)$.

Next fix opposite placements $1,2$.  Operators on distinct placements
commute, even when the players' local spaces differ.  Put
$W_b=Z_{1,b}Z_{2,b}$, $D_a=X_{1,a}-X_{2,a}$, and
$K_{j,a,b}=[X_{j,a},Z_{j,b}]$.  Define the nonnegative scalar errors
\[
 \begin{aligned}
 e_X&=\sum_a\|D_a\hat\psi\|^2,&
 e_Z&=\sum_b\|(Z_{1,b}-Z_{2,b})\hat\psi\|^2,\\
 k_j&=\sum_{a,b}\|K_{j,a,b}\hat\psi\|^2.&
 \end{aligned}
\]
For any two complete projective measurements $P,Q$ and any vector $w$,
expanding the square gives
\[
 \sum_c\operatorname{Re}\langle w,P_cQ_cw\rangle
   =\|w\|^2-\tfrac12\sum_c\|(P_c-Q_c)w\|^2.
\]
Cross-placement commutation gives
$R_{1,a,b}R_{2,a,b}=W_b^\dagger X_{1,a}X_{2,a}W_b$.
Applying the preceding identity first to the $X$ measurements in $W_b\hat\psi$
and then to the $Z$ measurements in $\hat\psi$ proves the exact defect formula
\[
 1-\sum_{a,b}\langle\hat\psi,R_{1,a,b}R_{2,a,b}\hat\psi\rangle
   =\tfrac12 e_Z+\tfrac12\sum_{a,b}\|D_aW_b\hat\psi\|^2.
\]
The overlap is real because the two placed sandwich effects commute.
To bound the last sum, use the algebraic identity
\[
 D_aW_b=W_bD_a+Z_{2,b}K_{1,a,b}-Z_{1,b}K_{2,a,b}.
\]
The squared norm of a sum of three vectors is at most three times the sum
of their squared norms.  Projector contraction and
$\sum_b\|W_bv\|^2\leq\sum_b\|Z_{2,b}v\|^2=\|v\|^2$ then bound the defect by
$c_{x,z}=\tfrac12 e_Z+\tfrac32(e_X+k_1+k_2)$.
Self-consistency of the point measurements gives
$\E e_X,\E e_Z=O(\eps)$; Parseval gives $\E k_j=O(\sqrt\eps)$.
Consequently $\E c_{x,z}\leq C(\eps+\sqrt\eps)$ for a universal $C$.\footnote{
The exact identity and estimate are \path{sandwich_overlap_identity} and
\path{sandwich_defect_pointwise_le} in
\path{MIPStarRE/QPBT/Combining/Points/Consistency.lean}.
They replace the cited Cauchy--Schwarz chain, not its conclusion.}

\section{Rounding and register transfer}

The orthonormalization lemma for two consistent bipartite POVMs returns a
projective measurement on the same local space as the first POVM.
Its explicit version bounds squared distance to that POVM by $220c^{1/4}$
when the consistency defect is at most $c$.  Apply it to the two sandwiches
for each $(x,z)$, separately for Alice and Bob.  Concavity of the fourth root
gives an averaged squared distance at most
$220(\E c_{x,z})^{1/4}$.\footnote{This uses
\path{projective_rounding_with_explicit_constant} in
\path{MIPStarRE/QPBT/Games/DistanceTheorems/ProjectiveRounding.lean}, the
explicit form of \path{lem:ortho}.  The fiberwise application and averaging
are in \path{MIPStarRE/QPBT/Combining/Points/Orthonormalization.lean}.}

The involution $\tau=(A'\ B'')(A''\ B')$ fixes $A,B$ and preserves
$\hat\psi$: it exchanges the two identical EPR pairs and reverses each pair.
It carries $AA'$ to $AB''$ and $BA''$ to $BB'$, so the same locally chosen
measurements have the same rounding errors on the second placement scheme.
No symmetry of the original strategy is needed.  This is not the blueprint's
permutation product $\sigma\theta$: that product sends $B'$ to $A'$, whereas
$\tau$ sends $B'$ to $A''$.\footnote{Compare
\path{lem:symmetric-equivalents-transfer} in
\path{blueprint/src/chapter/ch14_qpbt_observables.tex} with
\path{ProjectiveSetting.eprCrossSwap} in
\path{MIPStarRE/QPBT/Combining/Points/Placement.lean}.}

For each directed opposite pair, transport the projective factors one at a
time to obtain $d^2((ZX)_p,(XZ)_{p'})=O(\eps)$.
Combining this estimate with rounding, $d^2(R,ZX)=O(\sqrt\eps)$,
and commutation proves all three target inequalities by the squared-distance
triangle bound.  For $0\leq\eps\leq1$, their errors are bounded by
$K\eps^{1/8}$.  For $\eps>1$, all the operator families involved satisfy
$\sum_{a,b}F_{a,b}^\dagger F_{a,b}\leq I$, so their squared distances on a
unit state are at most $4$; enlarging $K$ covers this case as well.

\section{Conclusion}

The blueprint retains the source's combined-point statement and its
binary-refinement proof.  The formal construction proves that statement by
the direct argument above, with both ordered products, all directed opposite
placements, and the original expanded local spaces.\footnote{The blueprint
entry is \path{lem:qld-4-10} in
\path{blueprint/src/chapter/ch15_qpbt_combining.tex}; its construction theorem
is \leanid{MIPStarRE.QPBT.exists\_combinedPointsWitness}.}
There is no discrepancy in the combined-point conclusion under the existing
parameter and polynomial-error conventions, and no additional construction
is assumed.  The finite common-ancilla theorem is proved independently;
this direct construction does not use it or assume zero-state padding of
the fixed strategy.

\bibliographystyle{alpha}
\bibliography{references}

\end{document}
